\chapter{LLM MCP Integrations (Gemini / Claude / Copilot / OpenAI)}
\ChapterMeta{Scope:}{}
{Unified MCP-first integration strategy across major LLM clients.}

\section{Principle: one backend, many clients}

All clients should share the same implementation backend:

\begin{lstlisting}[language=bash]
python3 tools/run_mcp_server.py
\end{lstlisting}

Core exposed tools:
\begin{itemize}
  \item \cmd{doctor}
  \item \cmd{generate\_config}
  \item \cmd{review\_config}
  \item \cmd{validate\_predictions}
  \item \cmd{validate\_dataset}
  \item \cmd{eval\_coco}
  \item \cmd{run\_scenarios}
  \item \cmd{convert\_dataset} (optional)
\end{itemize}

Return payload policy (for all clients):
\begin{itemize}
  \item machine-readable JSON
  \item short textual summary in \cmd{summary}
  \item stable top-level keys: \cmd{ok}, \cmd{tool}, \cmd{summary}, \cmd{exit\_code}
  \item optional metadata: \cmd{meta.git\_sha}, \cmd{meta.manifest\_sha256}, runtime details
\end{itemize}

\section{Layer model (fixed)}

\begin{enumerate}
  \item Core layer: I/O envelope, auth/log hooks, error shaping.
  \item YOLOZU API layer: CLI/Python wrapper, argument normalization, manifest resolution.
  \item Job layer: long-running operations with \cmd{job\_id} tracking.
  \item Artifact layer: run artifacts and metadata return (git SHA / manifest hash / environment).
\end{enumerate}

\section{Quick start (MCP + Actions API)}

\subsection{MCP server}
\begin{lstlisting}[language=bash]
python3 tools/run_mcp_server.py
\end{lstlisting}

Official MCP Lite surface (deterministic, intended for AI-safe clients):
\begin{itemize}
  \item \cmd{doctor}
  \item \cmd{generate\_config}
  \item \cmd{review\_config}
\end{itemize}

Copy-paste sanity checks:
\begin{lstlisting}[language=bash]
python3 tools/run_mcp_server.py --print-tools > reports/mcp_tools.json
python3 tools/run_mcp_server.py --sample-generate-config > reports/ai_generate_config.json
python3 tools/run_mcp_server.py --sample-review-config reports/ai_generate_config.json > reports/ai_review_config.json
\end{lstlisting}

\subsection{Actions API (optional route)}
\begin{lstlisting}[language=bash]
python3 tools/run_actions_api.py --host 127.0.0.1 --port 8080 --workers 1
curl -sS http://127.0.0.1:8080/healthz
\end{lstlisting}

\subsection{Minimal end-to-end check}
\begin{lstlisting}[language=bash]
python3 tools/yolozu.py doctor --output reports/doctor.json
python3 tools/yolozu.py validate predictions data/smoke/predictions/predictions_dummy.json --strict
python3 tools/yolozu.py eval-coco --dataset data/smoke --predictions data/smoke/predictions/predictions_dummy.json --dry-run --output reports/mcp_eval_coco_dry.json
\end{lstlisting}

\section{Client routing}

\subsection{Gemini}
Use MCP first. Optionally use API tool-calling via OpenAPI endpoint.

\subsection{Claude}
Use MCP first with thin prompt wrappers; avoid duplicating tool logic.

\subsection{Copilot}
Use Copilot Extensions / VS Code participant that calls MCP backend (or API fallback).

\subsection{OpenAI}
Use MCP first; add GPT Actions when OpenAPI registration is required.

Optional API route:
\begin{lstlisting}[language=bash]
python3 tools/run_actions_api.py --host 127.0.0.1 --port 8080 --workers 1
# OpenAPI: http://127.0.0.1:8080/openapi.json
\end{lstlisting}

\subsection{Ollama (local LLM)}
Ollama runs models locally and can be used as the reasoning engine behind an MCP-capable client.
The \yolozu{} backend does not change: keep using the same \yolozu{} MCP server.

Start (or confirm) Ollama is running:
\begin{lstlisting}[language=bash]
ollama serve
ollama pull llama3.1
\end{lstlisting}

Then configure your client runtime to use Ollama as an \textbf{OpenAI-compatible} endpoint:
\begin{lstlisting}[language=bash]
# Typical values (client-dependent)
OPENAI_BASE_URL=http://127.0.0.1:11434/v1
OPENAI_API_KEY=ollama

# Example model identifiers are Ollama-local names
OPENAI_MODEL=llama3.1
\end{lstlisting}

Optional sanity checks:
\begin{lstlisting}[language=bash]
# Native Ollama API
curl -sS http://127.0.0.1:11434/api/tags

# OpenAI-compatible API
curl -sS http://127.0.0.1:11434/v1/models
\end{lstlisting}

Notes:
\begin{itemize}
  \item MCP tool execution still routes to \cmd{python3 tools/run\_mcp\_server.py}; only the LLM provider/model changes.
  \item Environment variable names are client-specific; some clients use \cmd{OPENAI\_API\_BASE} instead of \cmd{OPENAI\_BASE\_URL}.
  \item Model capability varies by local weights; if tool selection is unstable, prefer a stronger local model and keep prompts/tool wrappers thin.
  \item Verify the endpoint is reachable via \cmd{http://127.0.0.1:11434/v1/models} (OpenAI-compatible) in environments that support it.
\end{itemize}

\section{Operational recommendation}

Maintain feature parity by implementing behavior once in backend runner and exposing it through MCP/API adapters.

Operational guards in v2 layer implementation:
\begin{itemize}
  \item Allowlist is manifest-first (\path{tools/manifest.json}) with conservative fallback defaults.
  \item Path tokens are guarded (\cmd{..} rejected; absolute paths outside workspace rejected).
  \item Job state persistence to \cmd{runs/mcp\_jobs/*.json} for restart-safe status retrieval.
  \item Stale in-flight states (queued/running after restart) are surfaced as \cmd{unknown}.
  \item CLI execution timeout and output-size caps with explicit truncation metadata.
\end{itemize}

\section{Tool map (actual MCP/Actions surface)}

\subsection{Official support boundary (1.0.x)}
\begin{itemize}
  \item Guaranteed (deterministic/lightweight): \cmd{doctor}, \cmd{generate\_config}, \cmd{review\_config}, \cmd{validate\_predictions}
  \item Best-effort (environment dependent): long training jobs, TensorRT build/export, OpenCV CUDA/OpenVINO backends
\end{itemize}

\subsection{Synchronous tools}
\begin{itemize}
  \item Data/validation: \cmd{doctor}, \cmd{validate\_predictions}, \cmd{validate\_dataset}, \cmd{convert\_dataset}
  \item Inference/calibration: \cmd{predict\_images}, \cmd{parity\_check}, \cmd{calibrate\_predictions}
  \item Evaluation: \cmd{eval\_coco}, \cmd{eval\_instance\_seg}, \cmd{eval\_long\_tail}, \cmd{run\_scenarios}
\end{itemize}

\subsection{Asynchronous job tools}
\begin{itemize}
  \item \cmd{train\_job}, \cmd{export\_predictions\_job}, \cmd{test\_job}, \cmd{ttt\_job}, \cmd{ctta\_job}
  \item compatibility alias: \cmd{export\_onnx\_job} (same behavior as \cmd{export\_predictions\_job})
  \item control/query: \cmd{jobs.list}, \cmd{jobs.status}, \cmd{jobs.cancel}, \cmd{runs.list}, \cmd{runs.describe}
\end{itemize}

\section{AI operation pitfalls (high-impact)}

\begin{enumerate}
  \item Always pass \cmd{dataset} as a YOLO root and specify \cmd{split} explicitly.
  \item Prefer workspace-relative paths to avoid API-layer path guard rejections.
  \item Treat \cmd{ok/tool/summary/exit\_code} as the canonical status contract.
  \item For long tasks, do not poll stdout; use \cmd{job\_id} with \cmd{jobs.status}.
  \item Assume post-restart in-flight jobs become \cmd{unknown}; requeue deterministically if needed.
  \item Set \cmd{dry\_run}/\cmd{strict}/\cmd{force} explicitly to avoid client-default drift.
\end{enumerate}

\section{Efficiency backlog for MCP hardening}

High-value follow-ups that reduce AI ambiguity and integration drift:
\begin{itemize}
  \item Generated MCP/Actions reference now lives at:
  \begin{itemize}
    \item \path{docs/generated/mcp\_actions\_tool\_reference.json}
    \item \path{docs/generated/mcp\_actions\_tool\_reference.md}
  \end{itemize}
  \item Regenerate/check command: \cmd{python3 tools/generate\_integration\_tool\_reference.py --check}
  \item Add explicit error-code taxonomy for allowlist/path/timeout failures in integration responses.
\end{itemize}

\section{CI no-abandon policy}

For MCP-related changes:
\begin{enumerate}
  \item Run local validation (manifest + tests).
  \item Push and monitor latest CI result.
  \item If failed, patch immediately and rerun until green.
\end{enumerate}

\section{Cross references}

\begin{itemize}
  \item OpenAI details: \cmd{docs/openai\_mcp\_actions.md}
  \item Copilot details: \cmd{docs/copilot\_mcp\_integration.md}
  \item Shared overview: \cmd{docs/llm\_integrations.md}
\end{itemize}
